\documentclass[12pt,a4paper]{article}
\usepackage{amsmath}
\usepackage{amssymb}
\usepackage{bm}
\usepackage[utf8]{inputenc}
\usepackage{xeCJK}

\setCJKmainfont{Hiragino Sans GB}

\title{京都大学大学院 数理工学専攻}
\author{凸最適化 過去問}
\date{H27 (OR)}

\begin{document}

\maketitle

関数 $f:\mathbb{R}^n \to \mathbb{R}$ を連続的に微分可能な凸関数とし，
$S=\{x\in\mathbb{R}^n\mid a^T x=b\}$ とする。
ただし，$a$ は $0$ でない $n$ 次元ベクトル，$b$ はスカラー，$^T$ はベクトルの転置を表す。
次の凸計画問題を考える。

\[
\text{(P)}:\quad\text{Minimize}\quad f(x)\quad
\text{subject to}\quad x\in S
\]

さらにパラメータ $z\in\mathbb{R}^n$ を含む次の凸二次計画問題を考える。

\[
\text{P}(z):\quad\text{Minimize}\quad \nabla f(z)^T y + \frac{1}{2}(y-z)^T(y-z)\quad
\text{subject to}\quad y\in S
\]

ここで，決定変数は $y$ である。任意の $z\in\mathbb{R}^n$ に対して
問題 $\text{P}(z)$ は唯一の最適解 $\bar{y}(z)$ をもつ。
以下の問いに答えよ。

\section*{(i)}

$z\in S$ とする。問題 $\text{P}(z)$ のカルーシュ・キューン・タッカー（Karush--Kuhn--Tucker）条件を用いて
$\bar{y}(z)$ を求めよ。

\section*{(ii)}

$x\in S$ かつ $\bar{y}(x)=x$ であるとき，$x$ は問題 $\text{(P)}$ の最適解であることを示せ。

\section*{(iii)}

$x\in S$ かつ $\bar{y}(x)\neq x$ であるとき，
$\nabla f(x)^T(\bar{y}(x)-x) < 0$，
$a^T(\bar{y}(x)-x)=0$ であることを示せ。

\section*{(iv)}

$\bar{y}(x)\neq x$ であるとき，$x$ は問題 $\text{(P)}$ の最適解でないことを示せ。

\end{document}
